% =============================================================================
%  SLIDE 10 – Practice exercise (លំហាត់ប្រតិបត្តិ)          [2 slides]
%    10a  problem + figure (4 frames)       10b  proof (7 frames)
%         1: statement   2: figure   3: ∠ABC = ∠AED   4: AB = AE
% =============================================================================
\begin{frame}
\begin{tikzpicture}[remember picture,overlay,shift={(current page.south west)},x=1cm,y=1cm]
  \sHdr{លំហាត់ប្រតិបត្តិ}{Practice Exercise}

  \node[anchor=north west,text=NavyDeep,align=left,inner sep=0,font=\fontsize{12.5}{22}\selectfont] at (1.0,7.0) {ក្នុងរូបខាងស្តាំគេមាន $\angle ABC=\angle AED$\\
     និង $AB=AE$ ។ ចូរបង្ហាញថា\\
     \textcolor{RoyalBlue}{\textbf{ក.}} $\Delta ABC\cong\Delta AED$\\
     \textcolor{RoyalBlue}{\textbf{ខ.}} $AD=AC$ ។};
  \node[fill=GreenAccent!12,draw=GreenAccent!70!black,rounded corners=10pt,inner xsep=10pt,inner ysep=4pt,
        text=NavyDeep] at (4.2,1.5) {\fontsize{11}{14}\selectfont អាចធ្វើជាដៃគូ ឬជាក្រុមតូច};

  \tenfig{8.6}{3.9}{2.7}{5.4}
  \only<2->{\tenbase}
  \only<3->{\tAngBE{AngCol}{0.6}
    \node[anchor=west,text=AngCol,inner sep=0] at (9.3,7.2) {\fontsize{12}{15}\selectfont $\angle ABC=\angle AED$};}
  \only<4->{\tSides{SideCol}
    \node[anchor=west,text=SideCol,inner sep=0] at (9.3,6.6) {\fontsize{12}{15}\selectfont $AB=AE$};}
  \sFtr
\end{tikzpicture}
\end{frame}
